% NeurIPS 2024 — nips.sty is tailored to the project language at creation.

\documentclass{article}

\usepackage[T2A]{fontenc}
\usepackage[utf8]{inputenc}
\usepackage[russian,english]{babel}
\usepackage{amsmath}

\title{Мультиязычная детекция ИИ-текста: разработка архитектуры для существующих энкодеров}

\author{
  Герасимов~С.\,А.\\
}
\date{}
\begin{document}

\maketitle
\renewcommand{\abstractname}{Аннотация}
\begin{abstract}
% Аннотация (150–250 слов): кратко о проблеме, подходе и результатах.
Стремительное развитие больших языковых моделей (LLM) делает задачу обнаружения сгенерированного текста всё более актуальной. Однако большинство современных детекторов обучаются и тестируются преимущественно на английском языке, и их поведение в мультиязычной среде остаётся малоизученным. В данной работе мы испытали современный детектор GigaCheck \cite{gigacheck2026}(на основе модели Mistral-7B) на нескольких мультиязычных бенчмарках: SemEval-2024 Task 8 \cite{semeval2024}, IberAuTexTification \cite{iberautext2024}, DetectRL-X \cite{detectrlx2026} и BLUFF \cite{bluff2026}. Мы оцениваем качество детекции на 10+ языках, включая как популярные (английский, русский, китайский), так и низкоресурсные (йоруба и тд.).

Результаты показывают высокую нестабильность: на одних языках и датасетах GigaCheck демонстрирует F1-score, близкий к 1.0 (например, на греческом и испанском в BLUFF), что объясняется артефактами машинного перевода, а на других — катастрофически падает. Так, на новостном домене BLUFF (тексты GPT-4) F1-score для английского и русского языков составляет около 0.64, на новостных статьях модель полностью коллапсирует, предсказывая все тексты как сгенерированные ИИ (Recall для класса human = 0.0). Мы связываем это с различиями в стиле профессиональной журналистики и низким качеством многих датасетов, что является проблемой.

Работа носит практический характер и закладывает основу для будущих улучшений. Мы приходим к выводу, что текущие детекторы не являются языконезависимыми и не устойчивы к смене языка. Это мотивирует дальнейшие исследования в направлении архитектур с вспомогательной головой Language ID и мультизадачным обучением, способных выучить более глубокие, устойчивые признаки генерации, не привязанные к конкретному языку или стилю.
\end{abstract}

\newpage

\section{Введение}
Стремительное развитие больших языковых моделей (LLM), таких как GPT-4, Mistral и LLaMA, привело к тому, что сгенерированные тексты стали практически неотличимы от написанных человеком. Хотя эти технологии открывают огромные возможности для автоматизации, они также создают серьезные риски: от распространения дезинформации и спама до академического плагиата и подрыва доверия к цифровому контенту. В связи с этим задача автоматической детекции машинно-сгенерированного текста (AI-generated text detection) становится критически важной для поддержания информационной безопасности.

За последние годы исследовательское сообщество добилось значительного прогресса в решении этой задачи. Современные детекторы, основанные на дообучении трансформеров (например, mDeBERTa) или адаптации больших языковых моделей (например, GigaCheck на базе Mistral-7B \cite{gigacheck2026}), демонстрируют впечатляющие результаты. На популярных бенчмарках, таких как SemEval-2024 Task 8 \cite{semeval2024}, их F1-score достигает 0.95 и выше. однако у этого успеха есть обратная сторона: подавляющее большинство существующих моделей и датасетов ориентированы исключительно на английский язык. Также во многих датасетах тесты генерировались устаревшими моделями(gpt3 и тд.), поэтому получались такие большие значения метрик качества. Вопрос о том, насколько хорошо эти детекторы работают в условиях реального мультиязычного мира, остается открытым.

При переносе детекторов на другие языки исследователи сталкиваются с двумя фундаментальными проблемами. Модели, обученные на английском, плохо справляются с низкоресурсными языками, так как их токенизаторы не оптимизированы под эти системы письма. Во-вторых, это сдвиг домена (domain shift) и эффект машинного перевода (translationese). Когда ИИ-тексты генерируются на одном языке, а затем переводятся на другой, они оставляют характерные артефакты. Парадоксально, но иногда детекторы показывают на таких переведенных текстах идеальную точность (F1 = 1.0), не потому что они выучили признаки генерации, а потому что они находят «швы» от переводчика.

В данной работе мы проводим исследование современного мультиязычного детектора GigaCheck \cite{gigacheck2026}, чтобы понять его реальные возможности и ограничения. Мы тестируем модель на четырех крупных бенчмарках (SemEval-2024 \cite{semeval2024}, IberAuTexTification \cite{iberautext2024}, DetectRL-X \cite{detectrlx2026} и BLUFF \cite{bluff2026}), охватывающих более 10 языков — от популярных (английский, китайский, испанский) до редких (йоруба, малагасийский).

Наши эксперименты выявляют нестабильность. С одной стороны, GigaCheck показывает F1-score, близкий к 1.0, на греческом и испанском языках в датасете BLUFF, что объясняется артефактами перевода. С другой стороны, на родном английском языке в том же датасете F1 падает до 0.64, она предсказывает метку «ИИ» для всех текстов подряд (Recall для человеческих текстов падает до 0.0), принимая профессиональный журналистский стиль за машинную генерацию.

Этот результат наглядно демонстрирует, что текущие детекторы не являются языконезависимыми и не устойчивы к смене языка. Основная цель данной работы — предложить решение для мультиязычной детекции и проанализировать слабые места существующих детекторов. Мы утверждаем, что для создания надежного детектора необходимо отойти от моноязычных архитектур и использовать мультизадачное обучение (Multi-Task Learning) с добавлением вспомогательных голов, таких как определение языка (Language ID) \cite{advacheck2024, gigacheck_span2024}. Это позволит модели выучить более глубокие, универсальные признаки генерации, не привязываясь к конкретному языку или стилю.

\section{Постановка задачи}
\subsection{Общая постановка}
Мы решаем задачу бинарной классификации текста: определить, написан ли данный текст человеком (\texttt{human}) или сгенерирован языковой моделью (\texttt{ai}). 

Однако, в отличие от классических работ, мы фокусируемся не на английском языке, а на мультиязычном сценарии. Мы хотим понять, насколько современные детекторы (на примере GigaCheck) способны обобщаться на новые языки и домены, и предложить архитектуру, которая сделает детекцию более устойчивой к смене языка.

\subsection{Формальная постановка}

\textbf{Дано:}
\begin{itemize}
    \item Множество текстов \(X = \{x_1, x_2, \dots, x_N\}\).
    \item Каждый текст \(x_i\) принадлежит одному из языков \(L = \{l_1, l_2, \dots, l_K\}\) (например, 10 языков из BLUFF или DetectRL-X).
    \item Каждому тексту соответствует истинная метка \(y_i \in \{0, 1\}\), где 0 — человек, 1 — ИИ.
\end{itemize}

\textbf{Базовая задача (Single-Task):} \\
Обучить модель \(f_\theta\), которая минимизирует ошибку бинарной классификации:
\begin{equation}
    \min_\theta \sum_{i=1}^{N} \mathcal{L}_{cls}(f_\theta(x_i), y_i)
\end{equation}
где \(\mathcal{L}_{cls}\) — функция потерь (кросс-энтропия).

\textbf{Наша задача (Multi-Task):} \\
Мы хотим добавить вспомогательную задачу — определение языка. Тогда модель имеет две головы:
\begin{itemize}
    \item Голова классификации: \(f_{cls}(x_i) \rightarrow \{0, 1\}\).
    \item Голова языка: \(f_{lang}(x_i) \rightarrow \{1, \dots, K\}\).
\end{itemize}

Функция потерь становится суммой двух ошибок:
\begin{equation}
    \mathcal{L}_{total} = \mathcal{L}_{cls}(f_{cls}(x_i), y_i) + \alpha \cdot \mathcal{L}_{lang}(f_{lang}(x_i), l_i)
\end{equation}
где \(\alpha\) — коэффициент, балансирующий важность задач (например, \(\alpha = 0.5\)), а \(\mathcal{L}_{lang}\) — кросс-энтропия для определения языка.

\textbf{Гипотеза:} Обучение с дополнительной задачей Language ID заставит общий энкодер выучить такие представления текста, в которых языковые признаки отделены от признаков генерации. Это должно повысить F1-score на низкоресурсных языках.




\bibliographystyle{plainnat}
\begin{thebibliography}{99}

% --- Основные статьи, которые вы тестируете ---

\bibitem{gigacheck2026}
Tolstykh, I., Tsybina, A., Yakubson, S., Gordeev, A., Dokholyan, V., \& Kuprashevich, M. (2026).
\textit{GigaCheck: Detecting LLM-generated Content via Object-Centric Span Localization}.
Findings of the Association for Computational Linguistics: ACL 2026.

\bibitem{bluff2026}
Lucas, J., Murtagh-White, M., Uchendu, A., Al-Lawati, A., Yamashita, M., Macko, D., Srba, I., Moro, R., \& Lee, D. (2026).
\textit{BLUFF: Benchmarking in Low-resoUrce Languages for detecting Falsehoods and Fake news}.
Under Review 2026 --- Datasets and Benchmarks Track.

\bibitem{detectrlx2026}
Wang, Y., et al. (2026).
\textit{DetectRL-X: A Comprehensive Benchmark for Robust Evaluation of Machine-Generated Text Detectors}.
arXiv preprint arXiv:2605.15518.

% --- Статьи вашего научного руководителя (Андрей Грабовой) ---

\bibitem{advacheck2024}
Gritsai, G., Voznyuk, A., Khabutdinov, I., \& Grabovoy, A. (2024).
\textit{Advacheck at GenAI Detection Task 1: AI Detection Powered by Domain-Aware Multi-Tasking}.
arXiv preprint arXiv:2411.11736.

\bibitem{gigacheck_span2024}
Gritsai, G., Khabutdinov, I., \& Grabovoy, A. (2024).
\textit{Multi-head Span-based Detector for AI-generated Fragments in Scientific Papers}.
Proceedings of the Fourth Workshop on Scholarly Document Processing (SDP 2024), 220-225.

\bibitem{survey2024}
Gritsai, G., Voznyuk, A., Grabovoy, A., \& Chekhovich, Y. (2024).
\textit{Are AI Detectors Good Enough? A Survey on Quality of Datasets With Machine-Generated Texts}.
Preventing and Detecting LLM Misinformation (PDLM) at AAAI 2025. arXiv:2410.14677.

% --- Датасеты и соревнования ---

\bibitem{semeval2024}
Wang, Y., Mansurov, J., Ivanov, P., Su, J., Shelmanov, A., Tsvigun, A., ... \& Nakov, P. (2024).
\textit{SemEval-2024 Task 8: Multigenerator, Multidomain, and Multilingual Black-Box Machine-Generated Text Detection}.
Proceedings of the 18th International Workshop on Semantic Evaluation (SemEval-2024), 2024.

\bibitem{iberautext2024}
Sarvazyan, A. M., Ángel González, J., Rangel, F., Rosso, P., \& Franco-Salvador, M. (2024).
\textit{Overview of IberAuTexTification at IberLEF 2024: Detection and Attribution of Machine-Generated Text on Languages of the Iberian Peninsula}.
Procesamiento del Lenguaje Natural, 73(0), 421–434.


\end{thebibliography}

\end{document}